\chapter{Pseudorandomness}\label{ch:oc-pseudorandom}

This chapter introduces the structure-versus-randomness dichotomy, the Gowers norms that measure a failure of randomness, and the Green--Tao theorem as the application to the primes.

\section{The structure/randomness dichotomy}

A function on the integers is decomposed into a structured piece, which correlates with a nilsequence or a polynomial phase, and a pseudorandom piece, which is negligible in a chosen norm.

\section{Gowers norms}

The Gowers norm $\|f\|_{U^{k+1}}$ vanishes, in the limit, when $f$ has no correlation with polynomial phases of degree $k$. It is the norm in which the pseudorandom piece is small.

\section{The Green--Tao theorem}

The Green--Tao theorem states that the primes contain arithmetic progressions of every finite length. The proof transfers Szemer\'edi's theorem from the integers to a pseudorandom majorant of the primes.
